\subsection{The complete original relation rows in the graph endpoint control}

Keep exactly the domain and coefficient identity of (ACM17)--(ACM20).
The symbols in this paragraph include a superscript $\mathrm{rel}$
so that the monic relation polynomial is never confused with a
consecutive quotient-volume ratio. They are the literal objects of
(PGRT.2)--(PGRT.8):
\[
\begin{gathered}
 r_k(u)=\int_{\sum t_i=u}(1+|\tau_k|^2)\prod_iw_h(t_i),\qquad
 F_b=\operatorname{rem}_{h(s_1),\ldots,h(s_k)}(\tau_k\Sigma^b),\\
 z_b(u)=\int_{\sum t_i=u}\overline{\tau_k}F_b\prod_iw_h(t_i),
 \quad g_x=z(u)x/r_k(u),\quad
 F^0x=(sx)(k/2+iu)-g_x(u),\\
 d_j^{\rm rel}=\chi p_{j+1-q},\qquad
 a_j^{\rm rel}=\int|d_j^{\rm rel}(k/2+iu)|^2r_k(u)\,du>0,\\
 (\alpha_j^{\rm rel})_b=
 \int\overline{d_j^{\rm rel}(k/2+iu)}
       ((k/2+iu)^b r_k(u)-z_b(u))\,du,
 \quad j\geq q-1,\quad 0\leq b<q.
\end{gathered}\tag{ACM28}
\]
Here $s:E\to\mathcal P_{q-1}$ is the original monic remainder
section, $p_n$ is the unique monic polynomial of degree $n$
orthogonal to lower polynomials for $|\chi|^2r_k\,du$, and the
fibre integral eliminates $t_k$ with its original Lebesgue measure.
For $k=1$ it is evaluation. The constant factors are
$w_h(t)=|v_h(1/2+it)|^2/(2\pi)$ throughout. Each coefficient is
absolutely integrable by the proved polynomial exponential moments
and Cauchy--Schwarz applied to $F^0$. Expanding the full polynomials
in these integrals gives the finite products of the original
one-factor moments in (PGS.35) and (PGS.38), without dropping their
conjugate or cross terms. The coefficients of $p_n$ are obtained
from its positive finite lower-polynomial Gram system. This is
the exact finite arithmetic evaluation of each row.

The complete weighted density proof (PGRT.3)--(PGRT.6) gives
\[
 \widehat G_N-\Omega_k
  =\sum_{j=N}^{\infty}
       (a_j^{\rm rel})^{-1}(\alpha_j^{\rm rel})^*
                                  \alpha_j^{\rm rel},
 \qquad \Omega_k=K_{\rm low}+D_{\rm mix}^*D_{\rm mix}>0.
 \tag{ACM29}
\]
The series converges in the original finite matrix operator norm.
Specifically, the relation columns are a complete orthogonal family
in $L^2(r_kdu)$; the coefficient of $F^0x$ on its $j$th column is
$(a_j^{\rm rel})^{-1}\alpha_j^{\rm rel}x$. Removing the first
$N-q+1$ columns is precisely minimization over the original
relations $\chi\mathcal P_{N-q}$. Parseval gives the displayed
quadratic forms, and polarization gives their matrices. The trace
of the remaining positive tail is the sum of the squared tails
of the $q$ fixed frame vectors, hence tends to zero and bounds
its operator norm. The positive completed term and its mixed
realization are the same as (HC22), with full proof in PGMB.

For $q-1\leq j\leq2q-1$ put
\[
\begin{aligned}
 l_j^{\rm rel}
   &=(a_j^{\rm rel})^{-1}\alpha_j^{\rm rel}
             \widehat G_{q-1}^{-1}(\alpha_j^{\rm rel})^*,\\
 u_j^{\rm rel}
   &=(a_j^{\rm rel})^{-1}\alpha_j^{\rm rel}
             \widehat G_{2q}^{-1}(\alpha_j^{\rm rel})^*,\\
 Z_{2q}^{\rm mix}x&=(D_{\rm mix}x,f_{2q}x),\\
 u_j^{\rm rel}
   &=(a_j^{\rm rel})^{-1}\alpha_j^{\rm rel}
                     K_{\rm low}^{-1}(\alpha_j^{\rm rel})^*\\
   &\quad-(a_j^{\rm rel})^{-1}\left\|
 (I+Z_{2q}^{\rm mix}K_{\rm low}^{-1}(Z_{2q}^{\rm mix})^*)^{-1/2}
 Z_{2q}^{\rm mix}K_{\rm low}^{-1}(\alpha_j^{\rm rel})^*
                                      \right\|^2.
\end{aligned}\tag{ACM30}
\]
The domain of $Z_{2q}^{\rm mix}$ is the original $E$, and its
ordered target is the mixed weighted Hilbert space of (HC22)
followed by $L^2(r_kdu)$. In particular it retains the finite
tail as well as the mixed component. To prove the subtraction,
write $K=K_{\rm low}$ and $Z=Z_{2q}^{\rm mix}$. The identity
$\widehat G_{2q}=K+Z^*Z$ gives
\[
 (K+Z^*Z)^{-1}
 =K^{-1}-K^{-1}Z^*(I+ZK^{-1}Z^*)^{-1}ZK^{-1}.
\]
Multiplication by $K+Z^*Z$ cancels the leftover factor
$Z^*ZK^{-1}-Z^*(I+ZK^{-1}Z^*)(I+ZK^{-1}Z^*)^{-1}ZK^{-1}$.
The target inverse exists because it is the identity off the
finite-dimensional range of $Z$ and is positive with eigenvalues
at least one on that range. Pairing with the specified row proves
(ACM30). If $k\geq2$ and $\alpha_j^{\rm rel}\ne0$, the
subtracted squared norm is strictly positive: $D_{\rm mix}$ is
injective by the full fixed-sum calculation PGM.6--9, and the
displayed inverse square root has zero kernel. For $k=1$ its
mixed component is zero but its finite tail is retained. A zero
row gives zero in both terms. These statements retain the original
strictness domain and the whole negative term.

Use the endpoint weights $\omega_{q-1}=\omega_{2q-1}=1$ and
$\omega_j=2$ for $q\leq j\leq2q-2$, and define
\[
\begin{gathered}
 L_k^{\rm rel}=\sum_{j=q-1}^{2q-1}\omega_j\log(1+l_j^{\rm rel}),
 \qquad
 U_{k,2q}^{\rm rel}=\sum_{j=q-1}^{2q-1}\omega_j\log(1+u_j^{\rm rel}),\\
 L_k^{\rm rel}\leq\widehat{\mathcal B}_k
                                  \leq U_{k,2q}^{\rm rel},\\
 L_k^{\rm rel}-g_{s,\rm gr}^{+}\leq\mathcal B_k
                       \leq U_{k,2q}^{\rm rel}-g_{s,\rm gr}^{-}
 \qquad(s\geq0).
\end{gathered}\tag{ACM31}
\]
To verify every index, (ACM29) gives
$\widehat G_j=\widehat G_{j+1}
 +(a_j^{\rm rel})^{-1}(\alpha_j^{\rm rel})^*\alpha_j^{\rm rel}$.
The rank-one determinant identity therefore gives the $j$th
decrement $\log(1+(a_j^{\rm rel})^{-1}\alpha_j^{\rm rel}
\widehat G_{j+1}^{-1}(\alpha_j^{\rm rel})^*)$.
Since $\widehat G_{2q}\preceq\widehat G_{j+1}
\preceq\widehat G_{q-1}$, inverse order bounds its argument
between $1+l_j^{\rm rel}$ and $1+u_j^{\rm rel}$.
Inverse order follows by congruence with the smaller positive
form, inversion of its positive eigenvalues, and reversal of
that congruence. Summing the decrements from $q-1$ to $2q-2$
and from $q$ to $2q-1$ gives precisely the displayed weights.
The full quartet in (ACM17) has $q\geq4$, so the listed endpoints
and interior have these literal ranges. Finally (ACM20) says
$\mathcal B_k=\widehat{\mathcal B}_k-\Gamma_k^{\rm gr}$.
Subtracting (ACM27) yields the last line, with its ordered negative
$g_{s,\rm gr}^{-}$ term intact. This is a finite relation-row
bound on the same arithmetic scalar; it does not replace the
exact value of $\widehat{\mathcal B}_k$ when that value is retained.

